% Options for packages loaded elsewhere
% Options for packages loaded elsewhere
\PassOptionsToPackage{unicode}{hyperref}
\PassOptionsToPackage{hyphens}{url}
\PassOptionsToPackage{dvipsnames,svgnames,x11names}{xcolor}
\PassOptionsToPackage{space}{xeCJK}
%
\documentclass[
  11pt,
  letterpaper,
  oneside,
  a4paper]{book}
\usepackage{xcolor}
\usepackage[top=25.4mm,bottom=25.4mm,left=20mm,right=20mm,headheight=2.17cm,headsep=4mm,footskip=12mm]{geometry}
\usepackage{amsmath,amssymb}
\setcounter{secnumdepth}{5}
\usepackage{iftex}
\ifPDFTeX
  \usepackage[T1]{fontenc}
  \usepackage[utf8]{inputenc}
  \usepackage{textcomp} % provide euro and other symbols
\else % if luatex or xetex
  \usepackage{unicode-math} % this also loads fontspec
  \defaultfontfeatures{Scale=MatchLowercase}
  \defaultfontfeatures[\rmfamily]{Ligatures=TeX,Scale=1}
\fi
\usepackage{lmodern}
\ifPDFTeX\else
  % xetex/luatex font selection
  \ifXeTeX
    \usepackage{xeCJK}
    \setCJKmainfont[]{Songti SC}
    \setCJKsansfont[]{Heiti SC}
    \setCJKmonofont[]{FangSong\_GB2312}
  \fi
  \ifLuaTeX
    \usepackage[]{luatexja-fontspec}
    \setmainjfont[]{Songti SC}
  \fi
\fi
% Use upquote if available, for straight quotes in verbatim environments
\IfFileExists{upquote.sty}{\usepackage{upquote}}{}
\IfFileExists{microtype.sty}{% use microtype if available
  \usepackage[]{microtype}
  \UseMicrotypeSet[protrusion]{basicmath} % disable protrusion for tt fonts
}{}
\usepackage{setspace}
\makeatletter
\@ifundefined{KOMAClassName}{% if non-KOMA class
  \IfFileExists{parskip.sty}{%
    \usepackage{parskip}
  }{% else
    \setlength{\parindent}{0pt}
    \setlength{\parskip}{6pt plus 2pt minus 1pt}}
}{% if KOMA class
  \KOMAoptions{parskip=half}}
\makeatother
% Make \paragraph and \subparagraph free-standing
\makeatletter
\ifx\paragraph\undefined\else
  \let\oldparagraph\paragraph
  \renewcommand{\paragraph}{
    \@ifstar
      \xxxParagraphStar
      \xxxParagraphNoStar
  }
  \newcommand{\xxxParagraphStar}[1]{\oldparagraph*{#1}\mbox{}}
  \newcommand{\xxxParagraphNoStar}[1]{\oldparagraph{#1}\mbox{}}
\fi
\ifx\subparagraph\undefined\else
  \let\oldsubparagraph\subparagraph
  \renewcommand{\subparagraph}{
    \@ifstar
      \xxxSubParagraphStar
      \xxxSubParagraphNoStar
  }
  \newcommand{\xxxSubParagraphStar}[1]{\oldsubparagraph*{#1}\mbox{}}
  \newcommand{\xxxSubParagraphNoStar}[1]{\oldsubparagraph{#1}\mbox{}}
\fi
\makeatother


\usepackage{longtable,booktabs,array}
\newcounter{none} % for unnumbered tables
\usepackage{calc} % for calculating minipage widths
% Correct order of tables after \paragraph or \subparagraph
\usepackage{etoolbox}
\makeatletter
\patchcmd\longtable{\par}{\if@noskipsec\mbox{}\fi\par}{}{}
\makeatother
% Allow footnotes in longtable head/foot
\IfFileExists{footnotehyper.sty}{\usepackage{footnotehyper}}{\usepackage{footnote}}
\makesavenoteenv{longtable}
\usepackage{graphicx}
\makeatletter
\newsavebox\pandoc@box
\newcommand*\pandocbounded[1]{% scales image to fit in text height/width
  \sbox\pandoc@box{#1}%
  \Gscale@div\@tempa{\textheight}{\dimexpr\ht\pandoc@box+\dp\pandoc@box\relax}%
  \Gscale@div\@tempb{\linewidth}{\wd\pandoc@box}%
  \ifdim\@tempb\p@<\@tempa\p@\let\@tempa\@tempb\fi% select the smaller of both
  \ifdim\@tempa\p@<\p@\scalebox{\@tempa}{\usebox\pandoc@box}%
  \else\usebox{\pandoc@box}%
  \fi%
}
% Set default figure placement to htbp
\def\fps@figure{htbp}
\makeatother





\setlength{\emergencystretch}{3em} % prevent overfull lines

\providecommand{\tightlist}{%
  \setlength{\itemsep}{0pt}\setlength{\parskip}{0pt}}



 
\usepackage[]{natbib}
\bibliographystyle{plainnat-doi}


%==============================================================================
% ElegantBook PDF Style - LaTeX Preamble
% Replicates the visual style of ElegantBook template for Quarto PDF output
%==============================================================================

%------------------------------------------------------------------------------
% 1. Package Loading
%------------------------------------------------------------------------------
\usepackage{geometry}
\usepackage[table]{xcolor}
\usepackage[center]{titlesec}
\usepackage{fancyhdr}
\usepackage[most]{tcolorbox}
\usepackage[shortlabels,inline]{enumitem}
\usepackage{tikz}
\usetikzlibrary{backgrounds,calc,shadows,positioning,fit}
\usepackage{bbding}
\usepackage{manfnt}
\usepackage{adforn}
\usepackage{multicol}
\usepackage{amsmath,amssymb,amsfonts}
\usepackage{amsthm}
\usepackage[labelfont={bf,color=structurecolor}]{caption}
\captionsetup[table]{skip=3pt}
\captionsetup[figure]{skip=3pt}
\usepackage{booktabs}
\usepackage{multirow}
\usepackage{orcidlink}
\usepackage{etoolbox}

%------------------------------------------------------------------------------
% 2. Color Definitions (Blue Theme - Default)
%------------------------------------------------------------------------------
\definecolor{structurecolor}{RGB}{60,113,183}
\definecolor{main}{RGB}{0,166,82}
\definecolor{second}{RGB}{255,134,24}
\definecolor{third}{RGB}{0,174,247}
\definecolor{winered}{rgb}{0.5,0,0}
\definecolor{coverlinecolor}{RGB}{255,134,24}

% Color aliases for reference
\definecolor{structure3}{RGB}{60,113,183}
\definecolor{main3}{RGB}{0,166,82}
\definecolor{second3}{RGB}{255,134,24}
\definecolor{third3}{RGB}{0,174,247}

%------------------------------------------------------------------------------
% 3. Page Layout Adjustments (supplement to geometry in _extension.yml)
%------------------------------------------------------------------------------
% Note: Main geometry is set in _extension.yml
% Additional adjustments can be made here if needed

%------------------------------------------------------------------------------
% 4. Chapter/Section Format (titlesec) — match elegantbook.cls exactly
%------------------------------------------------------------------------------
% Chapter format - hang style (default)
% Note: elegantbook.cls uses \xchaptertitle which is "第 \thechapter{} 章"
% For Chinese chapter titles, set in \AtBeginDocument
\titleformat{\chapter}[hang]
  {\bfseries}
  {\filcenter\LARGE\bfseries\color{structurecolor} 第 \thechapter{} 章\;}
  {1pt}
  {\LARGE\bfseries\color{structurecolor}\filcenter}
  []

% Section format
\titleformat{\section}[hang]
  {\bfseries}
  {\Large\bfseries\color{structurecolor}\thesection\enspace}
  {1pt}
  {\color{structurecolor}\Large\bfseries\filright}

% Subsection format
\titleformat{\subsection}[hang]
  {\bfseries}
  {\large\bfseries\color{structurecolor}\thesubsection\enspace}
  {1pt}
  {\color{structurecolor}\large\bfseries\filright}

% Subsubsection format
\titleformat{\subsubsection}[hang]
  {\bfseries}
  {\large\bfseries\color{structurecolor}\thesubsubsection\enspace}
  {1pt}
  {\color{structurecolor}\large\bfseries\filright}

% Spacing
\titlespacing{\chapter}{0pt}{-20pt}{1.3\baselineskip}

%------------------------------------------------------------------------------
% 5. Headers and Footers (fancyhdr) — match elegantbook.cls
%------------------------------------------------------------------------------
\pagestyle{fancy}
\fancyhf{}

% Footer: centered page number
\fancyfoot[c]{\color{structurecolor}\small\thepage}

% Header: centered (chapter/section title)
\fancyhead[C]{\color{structurecolor}\normalfont\normalsize\leftmark}

% Header rule
\renewcommand{\headrule}{\color{structurecolor}\hrule width\textwidth}
\renewcommand{\headrulewidth}{1pt}

% Plain page style (chapter opening pages) — centered page number at bottom
\fancypagestyle{plain}{
  \renewcommand{\headrulewidth}{0pt}
  \fancyhf{}
  \fancyfoot[C]{\color{structurecolor}\small\thepage}
  \renewcommand{\headrule}{}
}

% Section and chapter marks
\renewcommand{\sectionmark}[1]{\markright{\thesection\ #1}}
\renewcommand{\chaptermark}[1]{\markboth{第 \thechapter{} 章\ #1}{}}

%------------------------------------------------------------------------------
% 6. List Environments (enumitem + tikz) — match elegantbook.cls
%------------------------------------------------------------------------------
\setlist{nolistsep}

% Itemize custom bullets
\newcommand*{\eitemi}{\tikz[baseline]\draw[baseline, ball color=structurecolor,draw=none] circle (2pt);}
\newcommand*{\eitemii}{\tikz[baseline]\draw[baseline, fill=structurecolor,draw=none,circular drop shadow] circle (2pt);}
\newcommand*{\eitemiii}{\tikz[baseline]\draw[baseline, fill=structurecolor,draw=none] circle (2pt);}

\setlist[itemize,1]{label={\eitemi}}
\setlist[itemize,2]{label={\eitemii}}
\setlist[itemize,3]{label={\eitemiii}}

% Enumerate format
\setlist[enumerate,1]{label=\color{structurecolor}\arabic*.}
\setlist[enumerate,2]{label=\color{structurecolor}(\alph*).}
\setlist[enumerate,3]{label=\color{structurecolor}\Roman*.}
\setlist[enumerate,4]{label=\color{structurecolor}\Alph*.}

%------------------------------------------------------------------------------
% 7. Theorem Environments (tcolorbox - fancy mode, match elegantbook.cls)
%------------------------------------------------------------------------------
% We define the tcolorbox styles here (same as elegantbook.cls),
% but defer the actual \newenvironment definitions to \AtBeginDocument
% (after Quarto's amsthm \newtheorem) to avoid conflicts.

\tcbuselibrary{theorems,breakable}

% Common style for all theorem boxes — match elegantbook.cls exactly
% Note: cls uses \citshape; for Chinese we use \kaishu via \citshape
% We'll define \citshape if not already defined (ctex provides \kaishu)
\providecommand{\citshape}{\itshape}

\tcbset{
  common/.style={
    fontupper=\citshape,
    lower separated=false,
    coltitle=white,
    colback=gray!5,
    boxrule=0.5pt,
    fonttitle=\bfseries,
    enhanced,
    breakable,
    top=8pt,
    before skip=8pt,
    attach boxed title to top left={yshift=-0.11in,xshift=0.15in},
    boxed title style={boxrule=0pt,colframe=white,arc=0pt,outer arc=0pt},
    separator sign={.},
  }
}

% defstyle (definition-like environments) — match cls exactly
\tcbset{
  defstyle/.style={
    common,
    colframe=main,
    colback=main!5,
    colbacktitle=main,
    overlay unbroken and last={
      \node[anchor=south east, outer sep=0pt] at (\linewidth-width,0) 
      {\textcolor{main}{$\clubsuit$}};
    },
  }
}

% thmstyle (theorem-like environments) — match cls exactly
\tcbset{
  thmstyle/.style={
    common,
    colframe=second,
    colback=second!5,
    colbacktitle=second,
    overlay unbroken and last={
      \node[anchor=south east, outer sep=0pt] at (\linewidth-width,0) 
      {\textcolor{second}{$\heartsuit$}};
    },
  }
}

% prostyle (proposition-like environments) — match cls exactly
\tcbset{
  prostyle/.style={
    common,
    colframe=third,
    colback=third!5,
    colbacktitle=third,
    overlay unbroken and last={
      \node[anchor=south east, outer sep=0pt] at (\linewidth-width,0) 
      {\textcolor{third}{$\spadesuit$}};
    },
  }
}

%------------------------------------------------------------------------------
% 7b. Theorem environment definitions moved to before-body.tex
%     (to run after Quarto's amsthm \newtheorem, avoiding conflicts)
%------------------------------------------------------------------------------

%------------------------------------------------------------------------------
% 8. Other Environments (example, exercise, problem, note, proof, solution, etc.)
%    Match elegantbook.cls exactly
%------------------------------------------------------------------------------
% Example environment
\newcounter{example}[chapter]
\renewcommand{\theexample}{\thechapter.\arabic{example}}
\newenvironment{example}[1][]{
  \refstepcounter{example}
  \par\noindent\textbf{\color{main}示例 \theexample #1 }\rmfamily}{
  \par\ignorespacesafterend}

% Exercise environment
\newcounter{exercise}[chapter]
\renewcommand{\theexercise}{\thechapter.\arabic{exercise}}
\newenvironment{exercise}[1][]{
  \refstepcounter{exercise}
  \par\noindent\makebox[-3pt][r]{\scriptsize\color{red!90}\HandPencilLeft\quad}
  \textbf{\color{main}练习 \theexercise #1 }\rmfamily}{
  \par\ignorespacesafterend}

% Problem environment
\newcounter{problem}[chapter]
\renewcommand{\theproblem}{\thechapter.\arabic{problem}}
\newenvironment{problem}[1][]{
  \refstepcounter{problem}
  \par\noindent\textbf{\color{main}问题 \theproblem #1 }\rmfamily}{
  \par\ignorespacesafterend}

% Note environment
\newenvironment{note}{
  \par\noindent\makebox[-3pt][r]{\scriptsize\color{red!90}\textdbend\quad}
  \textbf{\color{second}注} \citshape}{\par}

% proof, solution, remark environments are defined in before-body.tex
% (to run after Quarto's amsthm definitions)

% Assumption environment
\newenvironment{assumption}{\par\noindent\textbf{\color{third}假设} \citshape}{\par}

% Conclusion environment
\newenvironment{conclusion}{\par\noindent\textbf{\color{third}结论} \citshape}{\par}

% Property environment
\newenvironment{property}{\par\noindent\textbf{\color{third}性质} \citshape}{\par}

% Custom environment (user-defined name)
\newenvironment{custom}[1]{\par\noindent\textbf{\color{third} #1} \citshape}{\par}

%------------------------------------------------------------------------------
% 9. Introduction and Problemset Environments — match elegantbook.cls
%------------------------------------------------------------------------------

% Introduction environment style
\tcbset{
  introductionsty/.style={
    enhanced,
    breakable,
    colback=structurecolor!10,
    colframe=structurecolor,
    fonttitle=\bfseries,
    colbacktitle=structurecolor,
    fontupper=\citshape,
    attach boxed title to top center={yshift=-3mm,yshifttext=-1mm},
    boxrule=0pt,
    toprule=0.5pt,
    bottomrule=0.5pt,
    top=8pt,
    before skip=8pt,
    sharp corners
  }
}

% Introduction environment
% Note: Pandoc generates \begin{itemize}\tightlist\item...\end{itemize}
% inside the div. We wrap with tcolorbox + multicols but let the
% Pandoc-generated itemize provide the list structure.
\newenvironment{introduction}[1][内容提要]{
  \begin{tcolorbox}[introductionsty,title={#1}]
    \begin{multicols}{2}
      \renewcommand{\labelitemi}{\textcolor{structurecolor}{\upshape\scriptsize\SquareShadowBottomRight}}}{
    \end{multicols}
  \end{tcolorbox}}

% Problemset environment
% Pandoc generates \begin{enumerate}\tightlist\item...\end{enumerate}
% inside the div. We just provide the decorative header.
\newenvironment{problemset}[1][习题]{
  \vspace*{10pt}
  \begin{center}
    \textcolor{structurecolor}{\Large\bfseries\adftripleflourishleft~#1~\adftripleflourishright}
  \end{center}}{
  \par}

%------------------------------------------------------------------------------
% 10. Hyperref Settings (supplement to Quarto's default)
%------------------------------------------------------------------------------
% Note: colorlinks, linkcolor, citecolor, urlcolor are set in _extension.yml
% Additional hyperref settings

%------------------------------------------------------------------------------
% 11. Chinese Font Settings (for Chinese documents) — match elegantbook.cls ctexfont option
%------------------------------------------------------------------------------
% Requires XeLaTeX engine
% Use ctex with default fontset (auto-detect system fonts)
\usepackage[UTF8,scheme=plain]{ctex}
\xeCJKsetup{AutoFakeBold=true}
% Set Chinese fonts for macOS (adjust for other platforms)
%\setCJKmainfont{Songti SC}
%\setCJKsansfont{Heiti SC}
%\setCJKmonofont[AutoFakeBold=true]{FangSong_GB2312}

% Define \citshape (楷体 = \kaishu) if ctex is loaded
\ifdefined\kaishu
  \renewcommand{\citshape}{\kaishu}
\else
  \renewcommand{\citshape}{\itshape}
\fi
% Define \cfs (仿宋 = \fangsong) if ctex is loaded
\ifdefined\fangsong
  \newcommand{\cfs}{\fangsong}
\else
  \newcommand{\cfs}{\normalfont}
\fi

%------------------------------------------------------------------------------
% 12. Additional Styling Adjustments — match elegantbook.cls
%------------------------------------------------------------------------------

% Display math spacing
\setlength{\abovedisplayskip}{3pt}
\setlength{\belowdisplayskip}{3pt}

% Caption setup (already done above, repeated for clarity)
\captionsetup[table]{skip=3pt}
\captionsetup[figure]{skip=3pt}

% Listings style (code blocks) — match elegantbook.cls
\RequirePackage{listings}
\lstset{
  basicstyle=\ttfamily\small,
  breaklines=true,
  frame=single,
  rulecolor=\color{structurecolor},
  framerule=0.2pt,
  numbers=none,
  keywordstyle=\color{winered},
  commentstyle=\color{gray},
}

%------------------------------------------------------------------------------
% 13. Title Page — match elegantbook.cls \maketitle exactly
%------------------------------------------------------------------------------

% Chinese title page labels (from elegantbook.cls lang=cn)
\newcommand{\authorname}{\citshape 作者：}
\newcommand{\institutename}{\citshape 组织：}
\newcommand{\datename}{\citshape 时间：}
\newcommand{\versionname}{\citshape 版本：}

% Define ElegantBook-specific commands
\makeatletter
\newcommand{\subtitle}[1]{\gdef\@subtitle{#1}}
\newcommand{\institute}[1]{\gdef\@institute{#1}}
\newcommand{\version}[1]{\gdef\@version{#1}}
\newcommand{\extrainfo}[1]{\gdef\@extrainfo{#1}}
\newcommand{\logo}[1]{\gdef\@logo{#1}}
\newcommand{\cover}[1]{\gdef\@cover{#1}}
\newcommand\bioinfo[2]{\gdef\@bioinfo{{\citshape #1}：#2}}
\makeatother

% \cnormal is used in title page author/institute lines
\newcommand{\cnormal}{\kaishu}

% ---- Set cover and logo (defaults; override in project) ----
% These are set here so \maketitle can use them immediately.
% Users can override by redefining \@cover / \@logo before \maketitle.
\makeatletter
\gdef\@cover{images/cover.jpg}%
\gdef\@logo{images/logo-blue.png}%
\makeatother

% ---- Redefine \maketitle to match elegantbook.cls ----
% Must be in preamble (NOT \AtBeginDocument) so it overrides
% the standard \maketitle before \frontmatter\maketitle executes.
%
% The \@->\spacefactor error is caused by parskip.sty (loaded by
% Quarto's default template) which tries to process paragraph
% transitions after \includegraphics while still in vertical mode.
% Fix: wrap \includegraphics in \mbox{} to keep it in horizontal mode,
% then use explicit \par to transition back.
\makeatletter
\renewcommand*{\maketitle}{%
  \hypersetup{pageanchor=false}
  \pagenumbering{Alph}
  \begin{titlepage}
    \newgeometry{margin=0in}
    \parindent=0pt
    \parskip=0pt
    % Cover image — wrapped in \mbox to prevent
    % parskip \@ expansion in vertical mode
    \noindent\mbox{\includegraphics[width=\linewidth]{\@cover}}\par
    \nointerlineskip
    \setlength{\fboxsep}{0pt}%
    \noindent\colorbox{coverlinecolor}{\rule{0pt}{0.5in}\hspace{\linewidth}}\par
    \vfill
    \vskip-2ex
    \hspace{2em}\parbox{0.8\textwidth}{%
      \bfseries\Huge\@title\par}%
    \vfill
    \vspace{-1.0cm}
    \setstretch{2.5}
    \hspace{2.5em}%
    \begin{minipage}[c]{0.67\linewidth}
      {\color{darkgray}\bfseries\Large
        \ifcsname @subtitle\endcsname\@subtitle\\[2ex]\fi}
      \color{gray}\normalsize
      {\renewcommand{\arraystretch}{0.618}
      \begin{tabular}{l}
        \ifx\@author\empty\else\authorname\cnormal\@author\\ \fi
        \ifcsname @institute\endcsname \institutename \cnormal\@institute\\ \fi
        \datename\cnormal\zhtoday \\
        \ifcsname @version\endcsname \cnormal\versionname\@version\\ \fi
        \ifcsname @bioinfo\endcsname \cnormal\@bioinfo\\ \fi
      \end{tabular}}%
    \end{minipage}%
    \begin{minipage}[c]{0.27\linewidth}
    \begin{tikzpicture}[remember picture,overlay]
      \begin{pgfonlayer}{background}
        \node[opacity=0.8,
              anchor=south east,
              outer sep=0pt,
              inner sep=0pt] at ($(current page.south east) +(-0.8in,1.5in)$) {
                \ifcsname @logo\endcsname\includegraphics[width=4.2cm]{\@logo}\fi};
      \end{pgfonlayer}
    \end{tikzpicture}
    \end{minipage}
    \vfill
    \begin{center}
      \setstretch{1.3}
      \parbox[t]{0.7\textwidth}{\centering\citshape
        \ifcsname @extrainfo\endcsname\@extrainfo\fi}
    \end{center}
    \vfill
  \end{titlepage}
  \restoregeometry
  \thispagestyle{empty}%
  \pagenumbering{Roman}% switch to uppercase Roman numerals for TOC pages
}
\makeatother

%------------------------------------------------------------------------------
% 14. natbib citation style — numbered, sorted & compressed, square brackets
%------------------------------------------------------------------------------
\AtBeginDocument{%
  \setcitestyle{numbers,sort&compress,square}%
}

%------------------------------------------------------------------------------
% 15. Equation numbering — number within chapter (e.g. (1.1), (1.2), …)
%------------------------------------------------------------------------------
\numberwithin{equation}{chapter}

%------------------------------------------------------------------------------
% 16. Chapters start on odd pages; redefined \cleardoublepage to force
%     odd pages even with oneside (standard \cleardoublepage only does
%     \clearpage in oneside mode, which doesn't ensure odd page numbers).
%------------------------------------------------------------------------------
\makeatletter
\def\cleardoublepage{\clearpage\ifodd\c@page\else
    \hbox{}\thispagestyle{empty}\newpage\fi}
\makeatother

\AtBeginDocument{%
  \let\oldchapter\chapter
  \renewcommand{\chapter}{%
    \cleardoublepage
    \oldchapter
  }%
}%

%------------------------------------------------------------------------------
% 17. Equation reference parentheses — Quarto 对公式交叉引用使用 \ref
%     而非 \eqref，导致编号不带括号。AtBeginDocument 在 hyperref 加载后执行，
%     确保 \ref 重定义生效（覆盖 hyperref 版本）。
%------------------------------------------------------------------------------
\usepackage{xstring}
\AtBeginDocument{%
  \let\quartoOldRef\ref
  \renewcommand{\ref}[1]{%
    \IfBeginWith{#1}{eq-}{%
      \textup{(\quartoOldRef{#1})}%
    }{%
      \quartoOldRef{#1}%
    }%
  }%
}

%==============================================================================
% End of ElegantBook PDF Style Preamble (all other hooks cleaned up)
%==============================================================================
% Title page customization — ElegantBook style
\subtitle{Your Subtitle Here}
\institute{Your University, Your Department}
\version{v1.0}
\makeatletter
\@ifpackageloaded{bookmark}{}{\usepackage{bookmark}}
\makeatother
\makeatletter
\@ifpackageloaded{caption}{}{\usepackage{caption}}
\AtBeginDocument{%
\ifdefined\contentsname
  \renewcommand*\contentsname{Table of contents}
\else
  \newcommand\contentsname{Table of contents}
\fi
\ifdefined\listfigurename
  \renewcommand*\listfigurename{List of Figures}
\else
  \newcommand\listfigurename{List of Figures}
\fi
\ifdefined\listtablename
  \renewcommand*\listtablename{List of Tables}
\else
  \newcommand\listtablename{List of Tables}
\fi
\ifdefined\figurename
  \renewcommand*\figurename{Figure}
\else
  \newcommand\figurename{Figure}
\fi
\ifdefined\tablename
  \renewcommand*\tablename{Table}
\else
  \newcommand\tablename{Table}
\fi
}
\@ifpackageloaded{float}{}{\usepackage{float}}
\floatstyle{ruled}
\@ifundefined{c@chapter}{\newfloat{codelisting}{h}{lop}}{\newfloat{codelisting}{h}{lop}[chapter]}
\floatname{codelisting}{Listing}
\newcommand*\listoflistings{\listof{codelisting}{List of Listings}}
\usepackage{amsthm}
\theoremstyle{plain}
\newtheorem{corollary}{Corollary}[chapter]
\theoremstyle{plain}
\newtheorem{lemma}{Lemma}[chapter]
\theoremstyle{plain}
\newtheorem{proposition}{Proposition}[chapter]
\theoremstyle{plain}
\newtheorem{theorem}{Theorem}[chapter]
\theoremstyle{definition}
\newtheorem{definition}{Definition}[chapter]
\theoremstyle{remark}
\AtBeginDocument{\renewcommand*{\proofname}{Proof}}
\newtheorem*{remark}{Remark}
\newtheorem*{solution}{Solution}
\newtheorem{refremark}{Remark}[chapter]
\newtheorem{refsolution}{Solution}[chapter]
\makeatother
\makeatletter
\makeatother
\makeatletter
\@ifpackageloaded{caption}{}{\usepackage{caption}}
\@ifpackageloaded{subcaption}{}{\usepackage{subcaption}}
\makeatother
\makeatletter
\@ifpackageloaded{fontawesome5}{}{\usepackage{fontawesome5}}
\makeatother
\usepackage{bookmark}
\IfFileExists{xurl.sty}{\usepackage{xurl}}{} % add URL line breaks if available
\urlstyle{same}
\hypersetup{
  pdftitle={ElegantBook PDF Template},
  pdfauthor={Your Name},
  colorlinks=true,
  linkcolor={winered},
  filecolor={Maroon},
  citecolor={winered},
  urlcolor={winered},
  pdfcreator={LaTeX via pandoc}}


\title{ElegantBook PDF Template}


\author{\href{https://your-university.edu/your-profile}{Your
Name}\,~\orcidlink{0000-0000-0000-0001}}
\AtBeginDocument{\bioinfo{E-mail}{\texttt{your.email@university.edu}}}



\begin{document}
\frontmatter
\maketitle

%------------------------------------------------------------------------------
% Redefine theorem environments at \begin{document} time
% (after Quarto's amsthm \newtheorem, to avoid "already defined" error)
% Use \DeclareTColorBox to match elegantbook.cls visual style.
%------------------------------------------------------------------------------
% This runs in the normal token stream after \begin{document},
% so # does not need doubling.
%
% Quarto generates: \begin{definition}[name]\protect\hypertarget{id}{}\label{id}
% We accept the optional [name], and Quarto's extra tokens go into the box body.
%
% First, undefine Quarto's amsthm theorem environments
\makeatletter
\let\theorem\relax\let\endtheorem\relax
\let\definition\relax\let\enddefinition\relax
\let\lemma\relax\let\endlemma\relax
\let\corollary\relax\let\endcorollary\relax
\let\proposition\relax\let\endproposition\relax
\let\axiom\relax\let\endaxiom\relax
\let\postulate\relax\let\endpostulate\relax
\let\solution\relax\let\endsolution\relax
\let\remark\relax\let\endremark\relax
\makeatother

% Now define with simple \DeclareTColorBox (O{} for optional name)
% Note: Quarto puts \protect\hypertarget{}{}\label{} tokens BEFORE the theorem body.
% These end up inside the tcolorbox body and still function correctly.

% theorem (thmstyle, ♥) — 定理
\DeclareTColorBox[auto counter,number within=chapter]{theorem}{ O{} }{
  common,thmstyle,
  title={定理~\thetcbcounter\IfValueT{#1}{ (#1)}},
}

% definition (defstyle, ♣) — 定义
\DeclareTColorBox[auto counter,number within=chapter]{definition}{ O{} }{
  common,defstyle,
  title={定义~\thetcbcounter\IfValueT{#1}{ (#1)}},
}

% lemma (thmstyle, ♥) — 引理
\DeclareTColorBox[auto counter,number within=chapter]{lemma}{ O{} }{
  common,thmstyle,
  title={引理~\thetcbcounter\IfValueT{#1}{ (#1)}},
}

% corollary (thmstyle, ♥) — 推论
\DeclareTColorBox[auto counter,number within=chapter]{corollary}{ O{} }{
  common,thmstyle,
  title={推论~\thetcbcounter\IfValueT{#1}{ (#1)}},
}

% proposition (prostyle, ♠) — 命题
\DeclareTColorBox[auto counter,number within=chapter]{proposition}{ O{} }{
  common,prostyle,
  title={命题~\thetcbcounter\IfValueT{#1}{ (#1)}},
}

% axiom (thmstyle, ♥) — 公理
\DeclareTColorBox[auto counter,number within=chapter]{axiom}{ O{} }{
  common,thmstyle,
  title={公理~\thetcbcounter\IfValueT{#1}{ (#1)}},
}

% postulate (thmstyle, ♥) — 公设
\DeclareTColorBox[auto counter,number within=chapter]{postulate}{ O{} }{
  common,thmstyle,
  title={公设~\thetcbcounter\IfValueT{#1}{ (#1)}},
}

% Starred (unnumbered) versions
\DeclareTColorBox{theorem*}{ O{} }{
  common,thmstyle,
  title={定理\IfValueT{#1}{ (#1)}},
}
\DeclareTColorBox{definition*}{ O{} }{
  common,defstyle,
  title={定义\IfValueT{#1}{ (#1)}},
}
\DeclareTColorBox{lemma*}{ O{} }{
  common,thmstyle,
  title={引理\IfValueT{#1}{ (#1)}},
}
\DeclareTColorBox{corollary*}{ O{} }{
  common,thmstyle,
  title={推论\IfValueT{#1}{ (#1)}},
}
\DeclareTColorBox{proposition*}{ O{} }{
  common,prostyle,
  title={命题\IfValueT{#1}{ (#1)}},
}

% Redefine proof, solution, remark environments (match elegantbook.cls)
\makeatletter
\let\proof\relax\let\endproof\relax
\makeatother

% Proof environment — match elegantbook.cls
\newenvironment{proof}{
  \par\noindent\textbf{\color{second}\proofname\;}
  \color{black!90}\cfs}{\par}

% Solution environment — match elegantbook.cls
\newenvironment{solution}{\par\noindent\textbf{\color{main}解} \citshape}{\par}

% Remark environment — match elegantbook.cls
\newenvironment{remark}{\noindent\textbf{\color{second}注}}{\par}

% Chinese translation for proofname
\renewcommand{\proofname}{证明}

% Override Quarto's hardcoded "Table of contents"
% Quarto template unconditionally sets \contentsname to "Table of contents"
% just before \tableofcontents. We hook into \tableofcontents to re-override.
\pretocmd{\tableofcontents}{\renewcommand{\contentsname}{目\ \ 录}}{}{}

% Bibliography header suppression is now handled by before-bib.tex partial
% (natbib-based), which patches \bibliography to use \bibname and clears
% running headers. The old CSLReferences redefinition is no longer needed.

\renewcommand*\contentsname{Table of contents}
{
\setcounter{tocdepth}{1}
\tableofcontents
}

\setstretch{1.2}
\mainmatter
\pagenumbering{arabic}\pagestyle{fancy}

\bookmarksetup{startatroot}

\chapter*{Introduction}\label{introduction}
\addcontentsline{toc}{chapter}{Introduction}

\markboth{Introduction}{Introduction}

Introduction

\section*{Measure}\label{measure}
\addcontentsline{toc}{section}{Measure}

\markright{Measure}

\section*{Introduction}\label{introduction-1}
\addcontentsline{toc}{section}{Introduction}

\markright{Introduction}

\section*{IntroductionIntroductionIntroduction}\label{introductionintroductionintroduction}
\addcontentsline{toc}{section}{IntroductionIntroductionIntroduction}

\markright{IntroductionIntroductionIntroduction}

Introduction

\bookmarksetup{startatroot}

\chapter{ElegantBook PDF
扩展测试文档}\label{elegantbook-pdf-ux6269ux5c55ux6d4bux8bd5ux6587ux6863}

\begin{introduction}

\begin{itemize}
\tightlist
\item
  ElegantBook 样式测试
\item
  定理环境验证
\item
  列表环境测试
\item
  交叉引用功能
\end{itemize}

\end{introduction}

\section{简介}\label{ux7b80ux4ecb}

本文档用于测试 ElegantBook PDF 扩展的各项功能。

\section{定理环境测试}\label{ux5b9aux7406ux73afux5883ux6d4bux8bd5}

\subsection{定义环境}\label{ux5b9aux4e49ux73afux5883}

\begin{definition}[连续性]\protect\hypertarget{def-continuous}{}\label{def-continuous}

设函数
\(f: D \subseteq \mathbb{R} \to \mathbb{R}\)，\(x_0 \in D\)。如果对于任意
\(\epsilon > 0\)，存在 \(\delta > 0\)，使得当 \(|x - x_0| < \delta\) 且
\(x \in D\) 时，有 \(|f(x) - f(x_0)| < \epsilon\)，则称 \(f\) 在 \(x_0\)
处连续。

\end{definition}

\subsection{定理环境}\label{ux5b9aux7406ux73afux5883}

\begin{theorem}[中值定理]\protect\hypertarget{thm-mvt}{}\label{thm-mvt}

设函数 \(f: [a,b] \to \mathbb{R}\) 连续，在 \((a,b)\) 内可导。则存在
\(\xi \in (a,b)\)，使得

\begin{equation}\protect\phantomsection\label{eq-einstein}{
f'(\xi) = \frac{f(b) - f(a)}{b - a}
}\end{equation}

\end{theorem}

参考定理 Theorem~\ref{thm-mvt} 可知\ldots{}

\subsection{引理环境}\label{ux5f15ux7406ux73afux5883}

\begin{lemma}[费马引理]\protect\hypertarget{lem-fermat}{}\label{lem-fermat}

设函数 \(f\) 在 \(x_0\) 处可导，由公式 \ref{eq-einstein} 知 \(x_0\) 是
\(f\) 的极值点，则 \(f'(x_0) = 0\)。

\end{lemma}

\subsection{推论环境}\label{ux63a8ux8bbaux73afux5883}

\begin{corollary}[导数零点定理]\protect\hypertarget{cor-zero}{}\label{cor-zero}

若函数 \(f\) 在区间 \(I\) 上可导且 \(f'(x) = 0\) 对所有 \(x \in I\)
成立，则 \(f\) 在 \(I\) 上为常数函数。

\end{corollary}

\subsection{命题环境}\label{ux547dux9898ux73afux5883}

\begin{proposition}[极值必要条件]\protect\hypertarget{prp-extreme}{}\label{prp-extreme}

设函数 \(f\) 在 \(x_0\) 处可导，且 \(x_0\) 是 \(f\) 的局部极值点，则
\(f'(x_0) = 0\)。

\end{proposition}

\section{其他环境测试}\label{ux5176ux4ed6ux73afux5883ux6d4bux8bd5}

\subsection{示例环境}\label{ux793aux4f8bux73afux5883}

\begin{example}

考虑功能 \(f(x) = x^2\)。该函数在 \(x = 0\) 处可导，且 \(f'(0) = 0\)。

\end{example}

\subsection{练习环境}\label{ux7ec3ux4e60ux73afux5883}

\begin{exercise}

证明：若函数 \(f\) 在 \([a,b]\) 上连续，则 \(f\) 在 \([a,b]\) 上有界。

\end{exercise}

\subsection{问题环境}\label{ux95eeux9898ux73afux5883}

\begin{problem}

设 \(f(x) = \sin x\)，求 \(f'(x)\)。

\end{problem}

\subsection{注释环境}\label{ux6ce8ux91caux73afux5883}

\begin{note}

此处的证明需要用到拉格朗日中值定理。

\end{note}

\subsection{证明环境}\label{ux8bc1ux660eux73afux5883}

\begin{proof}
由拉格朗日中值定理，存在 \(\xi \in (a,b)\)，使得\ldots{}
\end{proof}

\subsection{解环境}\label{ux89e3ux73afux5883}

\begin{solution}
由 \(f'(x) = \cos x\)，得 \(f'(x) = \cos x\)。
\end{solution}

\subsection{备注环境}\label{ux5907ux6ce8ux73afux5883}

\begin{remark}
这个结果是微积分基本定理的直接推论。
\end{remark}

\subsection{假设环境}\label{ux5047ux8bbeux73afux5883}

\begin{assumption}

假设市场是无摩擦的，且不存在套利机会。

\end{assumption}

\subsection{结论环境}\label{ux7ed3ux8bbaux73afux5883}

\begin{conclusion}

综上所述，该期权定价公式在 Black-Scholes 模型下成立。

\end{conclusion}

\subsection{性质环境}\label{ux6027ux8d28ux73afux5883}

\begin{property}

布朗运动具有独立增量性和平稳增量性。

\end{property}

\section{列表环境测试}\label{ux5217ux8868ux73afux5883ux6d4bux8bd5}

\subsection{无序列表}\label{ux65e0ux5e8fux5217ux8868}

\begin{itemize}
\tightlist
\item
  第一项内容
\item
  第二项内容

  \begin{itemize}
  \tightlist
  \item
    嵌套第一项
  \item
    嵌套第二项

    \begin{itemize}
    \tightlist
    \item
      三级嵌套项
    \end{itemize}
  \end{itemize}
\end{itemize}

\subsection{有序列表}\label{ux6709ux5e8fux5217ux8868}

\begin{enumerate}
\def\labelenumi{\arabic{enumi}.}
\tightlist
\item
  第一步
\item
  第二步

  \begin{enumerate}
  \def\labelenumii{\alph{enumii}.}
  \tightlist
  \item
    子步骤 a
  \item
    子步骤 b

    \begin{enumerate}
    \def\labelenumiii{\roman{enumiii}.}
    \tightlist
    \item
      详细步骤 i
    \item
      详细步骤 ii
    \end{enumerate}
  \end{enumerate}
\end{enumerate}

\section{数学公式测试}\label{ux6570ux5b66ux516cux5f0fux6d4bux8bd5}

行内公式：\(E = mc^2\)

块级公式：

\begin{equation} \int_{-\infty}^{+\infty} e^{-x^2} dx = \sqrt{\pi} \end{equation}

\section{图表测试}\label{ux56feux8868ux6d4bux8bd5}

\subsection{表格}\label{ux8868ux683c}

\begin{longtable}[]{@{}lll@{}}
\caption{基本函数导数表 \{\#tbl:derivatives\}}\tabularnewline
\toprule\noalign{}
函数 & 导数 & 积分 \\
\midrule\noalign{}
\endfirsthead
\toprule\noalign{}
函数 & 导数 & 积分 \\
\midrule\noalign{}
\endhead
\bottomrule\noalign{}
\endlastfoot
\(x^n\) & \(nx^{n-1}\) & \(\frac{x^{n+1}}{n+1}\) \\
\(e^x\) & \(e^x\) & \(e^x\) \\
\(\ln x\) & \(\frac{1}{x}\) & \(x\ln x - x\) \\
\end{longtable}

\subsection{图片}\label{ux56feux7247}

\section{章后习题}\label{ux7ae0ux540eux4e60ux9898}

\begin{problemset}

\begin{enumerate}
\def\labelenumi{\arabic{enumi}.}
\tightlist
\item
  证明 Cauchy 收敛准则。
\item
  计算 \(\lim_{x \to 0} \frac{\sin x}{x}\)。
\item
  求解微分方程 \(y' = ky\)。
\end{enumerate}

\end{problemset}

\bookmarksetup{startatroot}

\chapter{第二章：高级主题}\label{ux7b2cux4e8cux7ae0ux9ad8ux7ea7ux4e3bux9898}

\section{随机控制理论初步}\label{ux968fux673aux63a7ux5236ux7406ux8bbaux521dux6b65}

随机控制是现代金融数学的重要工具。

\subsection{随机微分方程}\label{ux968fux673aux5faeux5206ux65b9ux7a0b}

考虑如下随机微分方程：

\begin{equation} dX_t = \mu(X_t, t) dt + \sigma(X_t, t) dW_t \end{equation}

其中 \(W_t\) 是标准布朗运动。

\begin{definition}[适应过程]\protect\hypertarget{def-adapted}{}\label{def-adapted}

随机过程 \(\{X_t\}_{t \geq 0}\) 称为适应的，如果对于每个
\(t \geq 0\)，\(X_t\) 是 \(\mathcal{F}_t\)-可测的。

\end{definition}

\subsection{HJB 方程}\label{hjb-ux65b9ux7a0b}

\begin{theorem}[HJB方程]\protect\hypertarget{thm-hjb}{}\label{thm-hjb}

值函数 \(V(x,t)\) 满足如下 HJB 方程：

\begin{equation} -\frac{\partial V}{\partial t} = \max_{u \in U} \left\{ \mathcal{L}^u V(x,t) + f(x,u,t) \right\} \end{equation}

其中 \(\mathcal{L}^u\) 是无穷小生成元。

\end{theorem}

\section{金融保险应用}\label{ux91d1ux878dux4fddux9669ux5e94ux7528}

\subsection{期权定价}\label{ux671fux6743ux5b9aux4ef7}

Black-Scholes 公式给出了欧式看涨期权的价格：

\begin{proposition}[Black-Scholes公式]\protect\hypertarget{prp-bs}{}\label{prp-bs}

欧式看涨期权的价格公式为：

\begin{equation} C(S,t) = S N(d_1) - K e^{-r(T-t)} N(d_2) \end{equation}

其中
\(d_1 = \frac{\ln(S/K) + (r + \sigma^2/2)(T-t)}{\sigma\sqrt{T-t}}\)，\(d_2 = d_1 - \sigma\sqrt{T-t}\)。

\end{proposition}

\subsection{最优投资消费}\label{ux6700ux4f18ux6295ux8d44ux6d88ux8d39}

\begin{problem}

在 CRRA 效用函数下，求解最优投资消费策略。

\end{problem}

\section{习题}\label{ux4e60ux9898}

\begin{problemset}

\begin{enumerate}
\def\labelenumi{\arabic{enumi}.}
\tightlist
\item
  证明 Itô 引理。
\item
  推导 Black-Scholes PDE。
\item
  求解 Merton 最优投资问题。
\end{enumerate}

\end{problemset}

\bookmarksetup{startatroot}

\chapter{参考文献}\label{ux53c2ux8003ux6587ux732e}

这里假设 Quarto
会自动生成参考文献页。如果需要自定义，可以在这里添加内容。

\bookmarksetup{startatroot}

\chapter{测试}\label{ux6d4bux8bd5}

\begin{theorem}[直线定理]\protect\hypertarget{thm-line}{}\label{thm-line}

直线方程可以写为： \[ y = mx + b \]

\end{theorem}

我们将通过三个步骤定义可测函数的积分。首先定义非负简单函数的积分。以下设
\(E\) 是 \(\mathcal{R}^n\) 中的可测集。

This file is a sample book demonstrating Quarto's capabilities with the
standard LaTeX \texttt{book} document class. It supports PDF (via LaTeX)
output. It is assumed that the reader has a basic working knowledge of
stochastic calculus \citet{karatzas1991brownian} and stochastic control
theory \citet{yong1999sch}.

\section{Introduction}\label{introduction-2}

Introduction \#\# acv

Introduction

\section{test for cosmo}\label{test-for-cosmo}

IntroductionIntroductionIntroduction

\bookmarksetup{startatroot}

\chapter{Test}\label{test}

这是什么

\section{This is chapter 2}\label{this-is-chapter-2}

这是一个测试

\bookmarksetup{startatroot}

\chapter*{参考文献}\label{ux53c2ux8003ux6587ux732e-1}
\addcontentsline{toc}{chapter}{参考文献}

\markboth{参考文献}{参考文献}

\pagestyle{plain}

\renewcommand{\bibsection}{}
\bibliography{references.bib}

% Suppress running headers on bibliography pages.
% Quarto's book template places this partial after \bibliography,
% so heading patches go here but they run too late for \bibliography
% itself. The heading is already set by Quarto's biblio-title.
\let\stdchaptermark\chaptermark
\renewcommand{\chaptermark}[1]{}
\makeatletter
\let\stdmkboth\@mkboth
\def\@mkboth#1#2{}
\makeatother
\markright{}
\pagestyle{plain}
\backmatter
\renewcommand\bibname{参考文献}


\pagestyle{empty}
\vspace*{18pt}
\noindent\textbf{Your Name}\par

Your bio text here. This will appear on the author page at the end of
the book.\par

\end{document}
